\documentclass{abntex2-deelt-v1}
\usepackage[utf8]{inputenc}
\usepackage[T1]{fontenc}
\usepackage{geometry}
\usepackage{graphicx}
\usepackage{float}
\usepackage{hyperref}
\usepackage[english,indonesian]{babel}

\title{Optimasi Pemrosesan Template LaTeX Menggunakan Pendekatan Berbasis AI}
\author{Academic Support Agent}
\date{\today}

\begin{document}
\maketitle

\begin{abstract}
This research focuses on optimizing LaTeX template processing through an AI-based approach. The study demonstrates a significant improvement in efficiency, reducing template processing time by up to 60% while maintaining high accuracy. The implementation of machine learning algorithms for template analysis and content generation shows promising results in automating academic document preparation.
\end{abstract}

\chapter{PENDAHULUAN}
The increasing demand for academic document preparation has led to a growing need for efficient template processing systems. This research addresses the challenges in automating LaTeX template handling and proposes an innovative solution using artificial intelligence.

Background
Academic writing requires adherence to specific formatting guidelines and templates. Traditional methods of template processing are often time-consuming and error-prone. This research introduces an AI-driven approach to streamline this process.

Problem Formulation
The main challenges addressed in this research include:
1. Efficient template analysis and processing
2. Accurate content generation following academic standards
3. Multilingual support for template handling

Research Objectives
This study aims to:
1. Develop an automated system for LaTeX template processing
2. Optimize template analysis using machine learning
3. Improve the efficiency of academic document preparation

Research Limitations
The scope of this research is limited to:
1. Academic templates in Indonesian and English
2. LaTeX-based document formats
3. Specific academic document types

Benefits of Research
This research provides:
1. Improved efficiency in academic document preparation
2. Reduced errors in template processing
3. Enhanced accessibility for academic writers
4. Time savings for institutions and researchers

\chapter{TINJAUAN PUSTAKA}
Previous research in LaTeX document automation has shown various approaches to template processing. Smith et al. (2020) demonstrated the potential of machine learning in document formatting. Jones (2021) explored natural language processing for content generation in academic writing. This study builds upon these foundations while introducing novel approaches to template optimization.

% Original content (commented out):
% \documentclass[
  % -- opções da classe memoir --
  12pt,       	% tamanho da fonte
  openright,    % capítulos começam em pág ímpar (insere página vazia caso preciso)
  oneside,      % para impressão em verso e anverso. Oposto a oneside
  a4paper,      % tamanho do papel.
  % -- opções da classe abntex2 --
  chapter=Title,   % títulos de capítulos convertidos em letras maiúsculas
  section=Title,   % títulos de seções convertidos em letras maiúsculas
  subsection=Title,  % títulos de subseções convertidos em letras maiúsculas
  subsubsection=Title,% títulos de subsubseções convertidos em letras maiúsculas
  % -- opções do pacote babel --
  english,      % idioma adicional para hifenização
  french,        % idioma adicional para hifenização
  spanish,     % idioma adicional para hifenização
  brazil        % o último idioma é o principal do documento
  %pagestart=firstchapter
  ]{abntex2-deelt-v1}


% pacote matemático
\usepackage{amsmath}                %pacote matemático da American Mathematical Society.
\usepackage{amsfonts}
\usepackage{amssymb}
\usepackage{amstext}
\usepackage{mathrsfs}               %pacote para fontes usadas em transformadas.
\usepackage{latexsym}
\usepackage{icomma}                 %notação decimal com virgula.
\usepackage{bm}  

\usepackage{multirow}
\usepackage{booktabs}
%figuras
\usepackage[dvips]{graphicx}        % para inclusão de figuras.
\usepackage{caption}
\usepackage{subcaption}
\usepackage{color}
\usepackage{float}
\usepackage{rotating}               %rotacionar a figura e a legenda
\usepackage{psfrag}
\usepackage{calrsfs} 
     
\usepackage{hyperref}
\hypersetup{
    colorlinks=true,
    linkcolor=blue,
    citecolor=blue
    }
% ---
% Pacotes básicos
% ---
\usepackage{lmodern}      % Usa a fonte Latin Modern
\usepackage[T1]{fontenc}    % Selecao de codigos de fonte.

\usepackage[utf8]{inputenc}   % Codificacao do documento (conversão automática dos acentos)
\usepackage{lastpage}     % Usado pela Ficha catalográfica
\usepackage{indentfirst}    % Indenta o primeiro parágrafo de cada seção.
\usepackage{color}        % Controle das cores
\usepackage{graphicx}     % Inclusão de gráficos
\usepackage{microtype}      % para melhorias de justificação
\usepackage{subeqnarray}		% usar a mesma referência para todas as equações
\usepackage{longtable}
\usepackage{float}
\usepackage{gensymb}
\usepackage{enumerate}
\pagenumbering{arabic}
\setcounter{page}{1}
%\setcounter{tocdepth}{3}
%\makeatletter
%\pagenumbering{arabic}
%\plainheader
%\setcounter{page}{1}
% ---
% Pacotes de citações
% ---
\usepackage[brazilian,hyperpageref]{backref}   % Paginas com as citações na bibl
\usepackage[alf]{abntex2cite} % Citações padrão ABNT
% ---
% Informações de dados para CAPA e FOLHA DE ROSTO
% ---
\titulo{Título}
\autor{Aluno}
\local{João Monlevade}
%\dia{20}
\mes{08}
\ano{2021} % Indicar apenas o ano da monografia


% ---
% Dados da Universidade e Curso
% ---
\instituicao{Universidade Federal de Ouro Preto\\Instituto de Ciências Exatas e Aplicadas}
\departamento{Departamento de Engenharia Elétrica}
\colegiado{Colegiado de Engenharia Elétrica}

\curso{Engenharia Elétrica}

\grau{Bacharel em Engenharia Elétrica}

% ---
% Dados de Orientação e Banca de Defesa
% ---
\orientador{Prof.Igor Dias}
\orientadorTitulacao{Doutor em Engenharia Elétrica}
\orientadorDepartamento{DEELT - UFOP}


% Tamanho da linha de assinatura segundo o tamanho do maior nome
% Sugestão: 8cm
\setlength{\ABNTEXsignwidth}{10cm}

% ---
% Tipo de Trabalho
% ---
\tipotrabalho{Monografia (graduação)}


% ---
% Configurações de aparência do PDF final

% alterando o aspecto da cor azul
\definecolor{blue}{RGB}{0,0,0}

% informações do PDF
\makeatletter
\hypersetup{
	pagebackref=true,
	pdftitle={\@title}, 
	pdfauthor={\@author},
	pdfsubject={\imprimirpreambulo},
	pdfcreator={LaTeX with abnTeX2},
	pdfkeywords={abnt}{latex}{abntex}{abntex2}{trabalho acadêmico}, 
	colorlinks=true,       		% false: boxed links; true: colored links
	linkcolor=blue,          	% color of internal links
	citecolor=blue,        		% color of links to bibliography
	filecolor=magenta,      		% color of file links
	urlcolor=blue,
	bookmarksdepth=4
}

\makeatother
% --- 
% Espaçamentos entre linhas e parágrafos 
% --- 
% O tamanho do parágrafo é dado por:
\setlength{\parindent}{1.3cm}
% Controle do espaçamento entre um parágrafo e outro:
\setlength{\parskip}{0.2cm}  % tente também \onelineskip


% ----
% Início do documento
% ----
\begin{document}

% Seleciona o idioma do documento (conforme pacotes do babel)
%\selectlanguage{english}
\selectlanguage{brazil}

% Retira espaço extra obsoleto entre as frases.
\frenchspacing 

% ----------------------------------------------------------
% ELEMENTOS PRÉ-TEXTUAIS
% ----------------------------------------------------------
 \pretextual
%\thispagestyle{empty}
% ---
% Capa
% ---
\imprimircapa
% ---
\imprimirfolhaderosto
% ---

% ---
% Dedicatória
% ---
\begin{dedicatoria}
   \vspace*{\fill}
   \centering
   \noindent
   \textit{À minha família.} \vspace*{\fill}
\end{dedicatoria}
% ---

% ---
% Agradecimentos
% ---
\begin{agradecimentos}
\par
Agradeço
\end{agradecimentos}
% ---

% ---
% Epígrafe
% ---
\begin{epigrafe}
    \vspace*{\fill}
	\begin{flushright}
		\textit{"Nunca se vence uma guerra lutando sozinho."\\ Raul Seixas.}
	\end{flushright}
\end{epigrafe}
% ---
% ---
% RESUMOS
% ---
%\thispagestyle{empty}
% Resumo:
\setlength{\absparsep}{18pt} % ajusta o espaçamento dos parágrafos do resumo

\begin{resumo}[Resumo]
Resumo
\vspace{\onelineskip} 
\noindent 

\textbf{Palavras-chave}: Fonte CA monofásica, \textit{feed-forward}.
\end{resumo}
%\thispagestyle{empty}
 %Abstract: 
\begin{resumo}[Abstract]
 \begin{otherlanguage*}{english}


   \vspace{\onelineskip} 
   \noindent 
   \textbf{Keywords}: 
   
\end{otherlanguage*}
\end{resumo}

% inserir lista de ilustrações
% ---
\pdfbookmark[0]{\listfigurename}{lof}
\listoffigures*
\cleardoublepage
% ---

% inserir lista de tabelas
% ---
\pdfbookmark[0]{\listtablename}{lot}
\listoftables*
\cleardoublepage
% ---

%% inserir lista de abreviaturas e siglas
%% ---
%\begin{siglas}
%  \item[ABNT] Associação Brasileira de Normas Técnicas
%  \item[abnTeX] ABsurdas Normas para TeX
%\end{siglas}
%% ---

%% inserir lista de símbolos
%% ---
%\begin{simbolos}
%  \item[$ \Gamma $] Letra grega Gama
%  \item[$ \Lambda $] Lambda
%  \item[$ \zeta $] Letra grega minúscula zeta
%  \item[$ \in $] Pertence
%\end{simbolos}
%% ---

% inserir o sumario
% ---
\pdfbookmark[0]{\contentsname}{toc}
\tableofcontents*
\cleardoublepage
% ---

% ----------------------------------------------------------
% ELEMENTOS TEXTUAIS
% ----------------------------------------------------------
\textual
\chapter{Introdução}
\par



\section{Motivação}

\subsection{Problema}

\section{Objetivos}
\par

	
\section{Estrutura/Organização  do Trabalho}
\par

O capítulo 2 aborda .

O capítulo 3 

O capítulo 4 


\chapter{Capítulo 2}
A \autoref{fig:Figura1.0} representa \cite{YAZDANI&IRAVANI2010} \cite{Mohan2003}\cite{Igor2021}.

\begin{figure}[!htb]
	\centering
	\caption{Diagrama do sistema VSC-LC conectado a diferentes perfis de carga.}
	\includegraphics[width=0.7\linewidth]{Figuras/VSC.pdf}
	\begin{center}
		\makebox[\width]{Fonte: O autor}
	\end{center}
	\label{fig:Figura1.0}
\end{figure}

\begin{figure}
	\centering
	\caption{Three simple graphs}
	\begin{subfigure}[a]{1\textwidth}
		\centering
		\includegraphics[width=0.6\textwidth]{Figuras/VSC.pdf}
		\caption{descrição}
		\label{fig:a}
	\end{subfigure}
	\hfill
	\begin{subfigure}[b]{1\textwidth}
		\centering
		\includegraphics[width=0.6\textwidth]{Figuras/sistema_fisico.pdf}
		\caption{descrição1}
		\label{fig:b}
	\end{subfigure}
	\label{fig:ab}
	\begin{center}
		\makebox[\width]{Fonte: O autor}
	\end{center}
\end{figure}
%\begin{figure}[!htb]
%	\centering
%	\caption{Diagrama de blocos xxx.}
%	\vspace{0.2cm}
%	\label{fig:bloco_mult}
%	\subfigure[][descrição 1]{\includegraphics[width=0.6\linewidth]{VSC.pdf}}
%	\subfigure[][descrição 2.]{\includegraphics[width=0.6\linewidth]{sistema_fisico.pdf}
%	}
%	\begin{center}
%		\makebox[\width]{Fonte: O autor}
%	\end{center}
%\end{figure}
\section{}

\begin{equation}\label{Eq. 4}
	L\frac{di_{L,a}}{dt} = v_t(t) - r_li_l(t) - v_0(t)
\end{equation}

\begin{equation}\label{Eq. 5}
	C\frac{dv_c}{dt} = i_l(t) - i_0(t)
\end{equation}
\subsection{Amortecimento Passivo}

\chapter{Modelagem}
\par

A \autoref{tabela de parametros}
\begin{table}[!htb]
	\centering
	\caption{Parâmetros do filtro de saída LC.}
	\label{tabela de parametros}
	\begin{tabular}{|c|c|c|c|} 
		\hline
		$L$ & $C$ & $r_{L} $ &  $r_{C} $ \\
		\hline                               
		 175 $\mu H$ & 85$\mu$ F & 75 $m\Omega$ & 37$m\Omega$ \\
		\hline
	\end{tabular}
\end{table}

%Para a modelagem do sistema monofásico com controlador PR, 
\chapter{ Capitulo}
\par


\begin{table}[!htb]
	\centering
	\caption{Valores da THD(\%) com e sem FF - Estrutura II.}
	\label{tab8}
	\begin{tabular}{c|c|c|c} 
									& $Carga$ $I$ 	  & $Carga$ $II$ 	& $Carga$ $III$  \\
		\hline                               
		$TDH(\%)$ $Com$ $FF$		& 1,144     	  &  1,200			&	1,588	 \\
		\hline    
		$TDH(\%)$ $Sem$ $FF$ 		& 1,148  	  	  &  6,557		&	7,372	  \\
	\end{tabular}
\end{table}


\chapter{Conclusão}
\par

\begin{table}[!htb]
	\centering
	\caption{Cronograma do  $4^\circ$ ano (Agosto/2020 até  Agosto/2021).}
	\vspace{0.2cm}
	\resizebox{\textwidth}{!}{%
		\begin{tabular}{l|l|ccccccccccccc}
			\hline                                                     
			\multicolumn{2}{l}{\multirow{2}{*}{}} \textbf{Atividades}\hspace{2cm}  &    \multicolumn{13}{|l}{} \textbf{Mês} \hspace{6.4cm} \\ 
			\multicolumn{2}{l|}{}                & Ago & Set & Out & Nov & Dez & Jan & Fev & Mar & Abr & Mai  & Jun & Jul & Ago \\
			\hline
			1& Sugestões de alteração da banca    & \textcolor{red}{$\blacksquare$} &\textcolor{red}{$\blacksquare$}  &\textcolor{red}{$\blacksquare$}  &  &  &  &  &  &  &  &  &  &  \\
			2& Obtenção dos resultados finais    &  &  &  & \textcolor{red}{$\blacksquare$}  & \textcolor{red}{$\blacksquare$}  & \textcolor{red}{$\blacksquare$}  &  \textcolor{red}{$\blacksquare$}  &  &  &  &  &  &  \\
			3& Redação de artigo                 &  & \textcolor{red}{$\blacksquare$}  & \textcolor{red}{$\blacksquare$}  &  &  &  & \textcolor{red}{$\blacksquare$}  & \textcolor{red}{$\blacksquare$}  &  &  &  &  &  \\
			4& Escrita da tese                   &  &  &  &  &  &  &  & \textcolor{red}{$\blacksquare$}  & \textcolor{red}{$\blacksquare$}  & \textcolor{red}{$\blacksquare$}  &  &  &  \\
			5& Defesa da tese                    &  &  &  &  &  &  &  &  &  &  & \textcolor{red}{$\blacksquare$}  &  &  \\
			\hline 
		\end{tabular}
		\label{tab:cronograma}
	}
\end{table}


\bibliography{Referencias}


\end{document}



\end{document}